\section*{Exercises}
\addcontentsline{toc}{section}{Exercises}

\subsection*{A. Theory}

\begin{exercise}
\exerciseid{13.A1}
State the spectral theorem for real symmetric matrices and explain why
orthogonality is part of the conclusion.
\end{exercise}

\begin{exercise}
\exerciseid{13.A2}
Let $A$ be symmetric. Prove that eigenvectors with distinct eigenvalues are
orthogonal.
\end{exercise}

\subsection*{B. By Hand}

\begin{exercise}
\exerciseid{13.B1}
Compute the Rayleigh quotient of
\[
  A=\begin{bmatrix}2&0\\0&5\end{bmatrix}
\]
at $x=(1,1)^T$ and at $x=(0,1)^T$.
\end{exercise}

\begin{exercise}
\exerciseid{13.B2}
For the same matrix $A$, find the minimum and maximum possible values of the
Rayleigh quotient.
\end{exercise}

\subsection*{C. Python / NumPy}

\begin{exercise}
\exerciseid{13.C1}
Sample random unit vectors and evaluate the Rayleigh quotient for a symmetric
$2\times 2$ matrix. Compare the observed minimum and maximum with
\texttt{np.linalg.eigvalsh}.
\end{exercise}

\begin{exercise}
\exerciseid{13.C2}
Implement 20 steps of power iteration for a symmetric matrix. Compare the final
Rayleigh quotient with the eigenvalue of largest absolute value.
\end{exercise}

\subsection*{D. Challenge}

\begin{exercise}
\exerciseid{13.D1}
Explain why the maximum Rayleigh quotient of a symmetric matrix should be an
eigenvalue. Use the orthonormal eigenbasis picture.
\end{exercise}
